% ============================ RESUMEN (TL;DR) ============================
\chapter*{Resumen}
\addcontentsline{toc}{chapter}{Resumen}

Se ha desplegado un laboratorio SOC sobre KVM en un único portátil, con
Wazuh 4.14 como SIEM, OPNsense como cortafuegos entre tres redes
aisladas y Suricata como IDS de red. Hay dos víctimas monitorizadas
(Windows 10 con Sysmon y Ubuntu Server) y una máquina atacante (Kali) en una red
separada, de modo que todo el tráfico del atacante atraviesa el cortafuegos.

Sobre este entorno se han ejecutado escenarios de ataque mapeados a MITRE
ATT\&CK y se ha documentado, para cada uno, la telemetría generada, la regla que
lo detecta, la cronología, la respuesta y los huecos de detección.

Posteriormente se modifican las reglas establecidas y se vuelven a poner a prueba con los mimos ataques. 

El objetivo es conseguir un sistema de monitoreo bien afinado que nos permita tener una visión clara de que sucede en nuestra infraestructura.

\begin{table}[H]\centering\small
\begin{tabularx}{\linewidth}{c X l l}
  \toprule
  \textbf{\#} & \textbf{Escenario} & \textbf{MITRE} & \textbf{Fase} \\
  \midrule
  1  & Escaneo de puertos con nmap (Kali $\rightarrow$ Windows y Ubuntu) & T1046     & Reconocimiento \\
  2  & Escaneo web con nikto (Kali $\rightarrow$ DVWA)                          & T1595.002 & Reconocimiento \\
  3  & Fuerza bruta de rutas con gobuster/dirb (Kali $\rightarrow$ DVWA)        & T1595.003 & Reconocimiento \\
  \addlinespace
  4  & Inyección SQL (Kali $\rightarrow$ DVWA)              & T1190     & Acceso inicial \\
  5  & XSS reflejado (navegador $\rightarrow$ DVWA)                             & T1059.007 & Acceso inicial \\
  6  & Subida de archivo malicioso (Kali $\rightarrow$ DVWA)                    & T1505.003 & Acceso inicial \\
  \addlinespace
  7  & Fuerza bruta SSH con hydra (Kali $\rightarrow$ Ubuntu)                   & T1110     & Acceso a credenciales \\
  8  & Fuerza bruta SSH con hydra (Kali $\rightarrow$ Windows)                  & T1110     & Acceso a credenciales \\
  9  & Fuerza bruta al login web con hydra http-form (Kali $\rightarrow$ DVWA)  & T1110     & Acceso a credenciales \\
  \addlinespace
  10 & Comandos en shell (Ubuntu)                                   & T1059.004 & Ejecución \\
  11 & PowerShell codificado con -EncodedCommand (Windows)                      & T1059.001 & Ejecución \\
  \addlinespace
  12 & Creación de usuario con useradd (Ubuntu)                                 & T1136.001 & Persistencia \\
  13 & Tarea programada con schtasks (Windows)                                  & T1053.005 & Persistencia \\
  14 & Clave de registro Run con reg add (Windows)                              & T1547.001 & Persistencia \\
  \addlinespace
  15 & Borrado del log de seguridad con wevtutil cl (Windows)                   & T1070.001 & Evasión \\
  \addlinespace
  16 & Reverse shell con netcat (Ubuntu $\rightarrow$ Kali)                     & T1059     & Command \& control \\
  \bottomrule
\end{tabularx}
\caption{Escenarios planificados}\label{tab:resumen}
\end{table}